
\setcounter{subsubsection}{3}

\begin{subproposition}
\label{XIII.4.7.4} Las notaciones son las de~\Ref{XIII.4.7}. Se supone
que, además de las condiciones \textup{a), b), c)}, se satisfacen las
condiciones siguientes:
\begin{enumerate}
\item[d)] Para todo punto $t$ de~$T=S-U$, el morfismo $f$ es localmente
$1$-asférico para~$\LL$ en $g(t)$.
\item[e)] Para todo punto $t\in T$, toda componente irreducible de la
fibra $X_t$ contiene a~$g(t)$ y, para todo punto $x$ de
$X_t-\{g(t)\}$ que no sea maximal, se tiene
$$
\prof\hop_x(X)\geq 3 \quad\textup{(SGA~2 XIV~1.2)}
$$
y el anillo $\cal{O}_{X,x}$ es noetheriano.
\end{enumerate}

Entonces el morfismo~\eqref{eq:XIII.4.7.3} es un isomorfismo.
\end{subproposition}

Como se ha dicho anteriormente, el grupo~$\pi$ se identifica con el
cociente de~$\pi'_1(X_U,a)$ por el subgrupo normal cerrado~$L$ engendrado
por la imagen de~$K$ \eqref{eq:XIII.4.7.2} en~$\pi'_1(X_U,a)$. Esto
equivale a decir que~$\pi$ clasifica los
\marginpar{425}
recubrimientos principales~$Z$ de~$X_U$ tales que $g^{-1}_U(Z)$ se
prolonga a un recubrimiento étale de~$S$, y que inducen sobre
$X_{\bar{s}}$ un recubrimiento obtenido por extensión del grupo
estructural a partir de un recubrimiento principal cuyo grupo es un
$\LL$-grupo. Para demostrar que~$\varphi$ es un isomorfismo, basta
demostrar que tal recubrimiento~$Z$ se prolonga a todo~$X$.

Demostremos primero que~$Z$ se prolonga a un recubrimiento étale de un
abierto que contiene a~$X_U$ y~$g(S)$. Sea~$W$ un esquema étale sobre~$X$
cuya imagen contiene a~$X_U$ y~$g(S)$, y pongamos $W_U=W\times_S U$. Del
hecho de que el morfismo $W\to S$ es $0$-acíclico y de la relación
$\mathrm{prof}.\,\et_T S\geq 2$ resulta que se tiene
$$
\mathrm{prof}.\,\et_{(W-W_U)} W\geq 2
$$
(SGA~2 XIV~1.13); por consiguiente, si $Z|W_U$ se prolonga a un
recubrimiento étale de~$W$, dicha prolongación es única salvo isomorfismo
único. De ello resulta que el problema de prolongar~$Z$ a un entorno de
$g(S)\cup X_U$ es local para la topología étale en un entorno de los
puntos de~$g(T)$. Si~$t$ es un punto de~$T$, pongamos $x=g(t)$ y denotemos
por~$\bar{X}$ (\resp por~$\bar{S}$) la localización estricta de~$X$ en
$\bar{x}$ (\resp de~$S$ en~$\bar{t}$), por
$\bar{U}=U\times_S\bar{S}$, por $\bar{X}_U=\bar{X}\times_X X_U$ y por
$\bar{g}\colon\bar{S}\to\bar{X}$ el morfismo deducido de~$g$. Basta
demostrar que, para todo punto~$t$ de~$T$, la imagen inversa~$\bar{Z}$
de~$Z$ sobre~$\bar{X}_U$ se prolonga a~$\bar{X}$ o, lo que equivale a lo
mismo, es trivial. Ahora bien, por definición de~$Z$, la imagen inversa
de~$\bar{Z}$ sobre~$\bar{U}$ es trivial. Para demostrar que~$\bar{Z}$ es
trivial, basta demostrar que es de la forma $\bar{f}^*_U E$, donde~$E$ es
un recubrimiento principal de~$\bar{U}$; en efecto, se tendrá
$E\simeq\bar{g}^*_U\bar{f}^*_U E\simeq\bar{g}^*_U\bar{Z}$, de donde se
sigue el resultado, puesto que $\bar{g}^*_U\bar{Z}$ es trivial. Ahora
bien, como el morfismo~$\bar{f}_U$ es $0$-acíclico y localmente
$0$-acíclico, para demostrar que~$\bar{Z}$ proviene de~$\bar{U}$ basta
demostrar que, para todo punto geométrico algebraico sobre un
punto de~$\bar{U}$, que puede suponerse que es el punto~$\bar{s}$,
\ifthenelse{\boolean{orig}}
{$\bar{Z}|X_{\bar{s}}$}
{$\bar{Z}|\bar{X}_{\bar{s}}$}
es trivial (SGA~4 XV~1.15). Pero, como
\ifthenelse{\boolean{orig}}
{$\bar{Z}|X_{\bar{s}}$}
{$\bar{Z}|\bar{X}_{\bar{s}}$}
se obtiene por extensión del grupo estructural a
partir de un recubrimiento principal cuyo grupo es un~$\LL$-grupo, esto
resulta del hecho de que el morfismo~$\bar{f}$ es $1$-asférico para~$\LL$.

Hemos
\marginpar{426}
demostrado, por tanto, que existe un entorno abierto~$V$ de
$g(S)\cup X_U$ tal que~$Z$ se prolonga a un recubrimiento étale~$Z_V$
de~$V$. Demostremos que~$Z_V$ se prolonga a todo~$X$. Basta ver que, para
todo punto~$x$ de~$X-V$, se tiene
$$
\prof\hop_x X\geq 3.
$$
Ahora bien, esto resulta de la hipótesis~e) y del hecho de que un punto
$x$ de~$X-V$ no puede ser maximal en su fibra~$X_t$, pues, como toda
componente irreducible de~$X_t$ contiene a~$g(t)$, todo punto maximal
de~$X_t$ pertenece a~$V$.
